\begin{table}[H]
\centering
\small
\caption{Incremental effects relative to M3 at the 10\% budget (254 cells). Every
Jaccard value and every ``shared'' count in this table uses the M3 top decile as
the reference. M4 is the full three-factor change, not the reference set.}
\label{tab:attribution_m3}
\begin{threeparttable}
\begin{tabularx}{\textwidth}{@{}>{\raggedright\arraybackslash}p{6.0cm} c c >{\raggedright\arraybackslash}X@{}}
\toprule
Scenario & Jaccard vs M3 & Shared / 254 & Interpretation \\
\midrule
M3 (reference) & 1.000 & 254 & M3 top decile \\
Reweighting only & 0.961 & 249 & Changes weights but keeps M3 components \\
Age swap only & 0.984 & 252 & Replaces total population with age count \\
Park addition only & 0.490 & 167 & Adds park term at matched weights \\
Age + reweighting & 0.954 & 248 & Joint change without park \\
Park + reweighting & 0.415 & 149 & Park dominates the joint change \\
Age + park & 0.499 & 169 & Park plus age substitution \\
M4 (all three changes) & 0.411 & 148 & Full M3-to-M4 construct change \\
\bottomrule
\end{tabularx}
\begin{tablenotes}[flushleft]\footnotesize
\item These are incremental comparisons against M3. They answer: how much does each change move the M3 top set? The corresponding M4-referenced comparisons are in Table~\ref{tab:attribution_m4}; the same numerical Jaccard cannot be interpreted as both an M3 and an M4 overlap. The exact three-factor Shapley shares are 85.46\% park, 13.25\% reweighting and 1.29\% age; these are descriptive shares of absolute rank movement, not causal effects.
\end{tablenotes}
\end{threeparttable}
\end{table}
